% ============================================================
%  Section: Inference Investigation
%  Paper: One Step Diffusion via Shortcut Models (Frans et al., ICLR 2025)
%  Capstone: Statistical Machine Learning
%  Author: Núi (tieunui)
% ============================================================

\section{Inference Investigation}
\label{sec:inference}

\subsection{Objectives}

The inference investigation is designed with three objectives:
(1)~validate the correctness of the Shortcut Model by comparing relative FID
across a range of step budgets (\textit{Number of Function Evaluations}, NFE);
(2)~confirm that \textit{shortcut conditioning} on parameter $d$ is the
essential factor governing sample quality, via an out-of-distribution (OOD)
ablation condition;
(3)~characterize the trade-off between inference cost and image quality within a
single unified model.

% ────────────────────────────────────────────────
\subsection{The FID Metric and Its Computation}
\label{sec:fid}

Fréchet Inception Distance (FID)~\cite{heusel2017gans} is the standard metric
for evaluating generative image models.
Rather than comparing images pixel-by-pixel, FID measures the distance between
two multivariate Gaussian distributions fitted to feature activations extracted
by Inception-v3~\cite{szegedy2016rethinking}, where the feature space has
dimension $D = 2048$.

\paragraph{Formula.}
Let $(\mu_r, \Sigma_r)$ denote the Inception feature statistics of the
reference image set (real data or a high-NFE proxy), and $(\mu_g, \Sigma_g)$
those of the generated set. FID is defined as:

\begin{equation}
    \mathrm{FID} = \|\mu_r - \mu_g\|^2
    + \mathrm{tr}\!\left(\Sigma_r + \Sigma_g
      - 2\left(\Sigma_r \Sigma_g\right)^{1/2}\right)
    \label{eq:fid}
\end{equation}

\noindent where $\mathrm{tr}(\cdot)$ denotes the matrix trace and
$(\Sigma_r \Sigma_g)^{1/2}$ is the matrix square root of the product.
A lower FID indicates that the generated distribution is closer to the
reference distribution.

\paragraph{Statistics structure.}
The file \texttt{celeba256\_fidstats\_ours.npz} provided by the authors stores
two arrays:
\begin{itemize}
    \item \texttt{mu}: mean feature vector of shape $(2048,)$,
          with $\bar{\mu} = 0.2774$;
    \item \texttt{sigma}: covariance matrix of shape $(2048, 2048)$,
          with $\bar{\sigma} = 2.11 \times 10^{-3}$.
\end{itemize}

The small value of $\bar{\sigma}$ indicates that Inception features on
CelebA-HQ are tightly clustered in $\mathbb{R}^{2048}$, consistent with the
highly structured nature of a face dataset.

\paragraph{Absolute vs.\ relative FID.}
The original paper reports \emph{absolute} FID-50k (generated samples vs.\
real CelebA-HQ 256 statistics, $n = 50\,000$).
The experiments in this report use \textbf{relative FID}: each NFE
configuration is compared against 128-step samples from the same checkpoint as
a proxy reference ($n = 504$ per run).
This is a practical choice given CPU-only compute constraints, and it measures
how much quality degrades as NFE decreases, rather than absolute sample
fidelity.

% ────────────────────────────────────────────────
\subsection{Experimental Setup}
\label{sec:setup}

\subsubsection{Environment and Libraries}

All experiments were run on a CPU-only macOS ARM (Apple Silicon) machine
without GPU or TPU acceleration.
Table~\ref{tab:env} lists the pinned library versions used to ensure
reproducibility.

\begin{table}[h]
\centering
\caption{Library versions used in the inference experiments.}
\label{tab:env}
\begin{tabular}{lll}
\toprule
\textbf{Library} & \textbf{Version} & \textbf{Notes} \\
\midrule
JAX / JAXlib     & 0.4.23 & Incompatible with 0.7+ when using Flax 0.7.4 \\
Flax             & 0.7.4  & \\
Optax            & 0.1.7  & \\
Orbax-checkpoint & 0.4.8  & v0.11+ requires JAX $\geq$ 0.6.0 \\
TensorFlow       & 2.21.0 & \texttt{tensorflow-cpu} unavailable for Python 3.12 ARM \\
pytorch-fid      & 0.3.0  & Relative FID computation, CPU-only \\
\bottomrule
\end{tabular}
\end{table}

\subsubsection{Model and Checkpoint}

The model used is \textbf{DiT-B}~\cite{peebles2023scalable}
(\textit{Diffusion Transformer, Base}) with the following hyperparameters:
\begin{itemize}
    \item \texttt{hidden\_size} = 768, \texttt{depth} = 12,
          \texttt{num\_heads} = 12, \texttt{patch\_size} = 2;
    \item Working space: latent $32 \times 32 \times 4$
          (not pixel space $256 \times 256 \times 3$);
    \item VAE decoder: \texttt{pcuenq/sd-vae-ft-mse-flax}
          (loaded from HuggingFace);
    \item Checkpoint: \texttt{celeba-shortcut2-every4400001}
          (training step 400\,000), downloaded from the authors'
          repository~\cite{Frans2025}.
\end{itemize}

\subsubsection{Sampling Procedure (Euler Sampling)}

The original repository requires a custom TFDS dataset builder
(\texttt{kvfrans/tfds\_builders}) not available in standard TFDS.
To avoid this dependency, a standalone inference script was developed that
samples Gaussian noise directly, bypassing the dataset loader.
The sampling loop follows Algorithm~2 of the paper~\cite{Frans2025}:

\begin{equation}
    x \leftarrow x + s_\theta(x,\, t,\, d)\cdot d, \qquad
    t \leftarrow t + d
    \label{eq:euler}
\end{equation}

\noindent where $d = 1/M$ and $M$ is the number of denoising steps.
The shortcut conditioning signal is encoded as:

\begin{equation}
    d_{\mathrm{flow}} = \lfloor \log_2 M \rfloor \in \mathbb{Z}_{\geq 0}
    \label{eq:dtflow}
\end{equation}

\noindent which reflects the binary-tree structure of the bootstrapping targets
used during training.

\subsubsection{Experiment Runs}

A total of six task groups were executed, covering 8 NFE configurations and
1 OOD ablation condition.
Each configuration generates $n = 504$ samples ($63 \times 8$ batches)
with \texttt{random\_seed = 42} unless otherwise specified.

\begin{table}[h]
\centering
\caption{Inference runs on CelebA-HQ 256 (DiT-B, training step 400\,000).}
\label{tab:runs}
\begin{tabular}{lccl}
\toprule
\textbf{Configuration} & \textbf{NFE} & $d_{\mathrm{flow}}$ & \textbf{Description} \\
\midrule
Shortcut (1-step)   & 1   & 0 & Single-step generation \\
Shortcut (2-step)   & 2   & 1 & \\
Shortcut (4-step)   & 4   & 2 & \\
Shortcut (8-step)   & 8   & 3 & \\
Shortcut (16-step)  & 16  & 4 & \\
Shortcut (32-step)  & 32  & 5 & \\
Shortcut (64-step)  & 64  & 6 & \\
Shortcut (128-step) & 128 & 7 & \textbf{Reference} \\
OOD baseline        & 128 & 0 (forced) & Shortcut model with conditioning $d=0$ \\
\bottomrule
\end{tabular}
\end{table}

% ────────────────────────────────────────────────
\subsection{Results}
\label{sec:results}

\subsubsection{Relative FID vs.\ NFE}

Table~\ref{tab:fid_relative} reports the relative FID measured against the
128-step reference samples.
Figure~\ref{fig:fid_nfe} illustrates the monotonically decreasing trend of FID
as a function of NFE on a log scale.

\begin{table}[h]
\centering
\caption{Relative FID by number of denoising steps (NFE) on CelebA-HQ 256.
Reference: 128-step samples, $n=504$. Lower FID is better.}
\label{tab:fid_relative}
\begin{tabular}{ccc}
\toprule
\textbf{NFE} & \textbf{Relative FID} & \textbf{Notes} \\
\midrule
1   & 0.2205 & \\
2   & 0.2007 & \\
4   & 0.1659 & $\dagger$ \\
8   & 0.1321 & \\
16  & 0.0925 & \\
32  & 0.0522 & \\
64  & 0.0182 & \\
128 & ---    & \textit{reference} \\
\midrule
OOD ($d=0$, 128 steps) & 100.51 & Out-of-distribution ablation \\
\bottomrule
\multicolumn{3}{l}{\footnotesize $\dagger$ NFE=4: FID = $0.2745 \pm 0.0811$
  (mean $\pm$ std over 3 seeds: 42, 123, 2025).}
\end{tabular}
\end{table}

\noindent The results show a monotonically decreasing FID as NFE increases,
with the steepest drop occurring between NFE=1 and NFE=8.
From NFE=16 onward, the curve flattens, exhibiting \textit{diminishing
returns}, consistent with theoretical expectations.

\subsubsection{Ablation: OOD Conditioning}

When $d$ is forced to $0$ at inference time (while the model was trained with
$d \in (0,1]$), the relative FID reaches $100.51$---approximately
$455\times$ higher than the properly conditioned 1-step result ($0.22$).
This confirms that shortcut conditioning is an indispensable component: once
the model receives inputs outside its training distribution, sample quality
collapses entirely.

\textbf{Interpretive note:} This OOD condition is \emph{not} equivalent to a
standalone Flow Matching model.
It represents a shortcut model operating under distribution shift, and the
result reflects a breakdown of the conditioning mechanism rather than a
direct architectural comparison between two separately trained models.

\subsubsection{Seed Stability}

At NFE=4, the experiment was repeated with three random seeds:

\begin{center}
\begin{tabular}{cc}
\toprule
\textbf{Seed} & \textbf{Relative FID} \\
\midrule
42   & 0.1659 \\
123  & 0.2969 \\
2025 & 0.3608 \\
\midrule
Mean $\pm$ Std & $0.2745 \pm 0.0811$ \\
\bottomrule
\end{tabular}
\end{center}

\noindent The standard deviation $\sigma = 0.081$ reflects non-negligible
variance given $n = 504$ and only 3 seeds.
This estimate should be interpreted as a rough stability indicator rather than
a statistically tight confidence interval, which would require at least 5 seeds
and bootstrap resampling.

\begin{figure}[h]
    \centering
    \includegraphics[width=0.72\textwidth]{figures/fid_nfe_curve.png}
    \caption{Relative FID vs.\ NFE (log base-2 scale) on CelebA-HQ 256 (DiT-B).
    The OOD point ($d=0$, FID $= 100.51$) is excluded from the axis
    to preserve readability.
    The curve confirms a monotonically decreasing trend with diminishing
    returns from NFE $\geq$ 16.}
    \label{fig:fid_nfe}
\end{figure}

\begin{figure}[h]
    \centering
    \includegraphics[width=0.85\textwidth]{figures/visual_grid_4x4.png}
    \caption{$4\times4$ image grid generated from the same initial noise vector
    at NFE $\in \{1, 4, 16, 128\}$.
    Observation: NFE=1 produces blurry images with visible texture artifacts;
    NFE=4 substantially improves facial structure; NFE=16 and NFE=128
    achieve visually near-equivalent quality.}
    \label{fig:visual_grid}
\end{figure}

% ────────────────────────────────────────────────
\subsection{Comparison with Authors' Results}
\label{sec:comparison}

Table~\ref{tab:comparison} places our experimental results alongside the
absolute FID-50k values reported by Frans et al.~\cite{Frans2025}.
It is important to note that the two metrics measure \emph{different
quantities} and cannot be compared numerically.

\begin{table}[h]
\centering
\caption{Qualitative FID trend comparison between this investigation and the
original results (Frans et al., ICLR 2025) on CelebA-HQ 256, DiT-B.}
\label{tab:comparison}
\begin{tabular}{lcccc}
\toprule
& \multicolumn{2}{c}{\textbf{This Work (Relative FID)}}
& \multicolumn{2}{c}{\textbf{Authors (FID-50k, Absolute)}} \\
\cmidrule(lr){2-3}\cmidrule(lr){4-5}
\textbf{NFE} & \textbf{Value} & \textbf{Trend}
             & \textbf{Value} & \textbf{Trend} \\
\midrule
1   & 0.2205 & \multirow{3}{*}{$\downarrow$ monotone} & 20.5 & \multirow{3}{*}{$\downarrow$ monotone} \\
4   & 0.1659 & & 13.8 & \\
128 & ---\ (ref.) & & 6.9  & \\
\midrule
OOD ($d=0$) & 100.51 & $\uparrow\uparrow$ & --- & --- \\
\bottomrule
\multicolumn{5}{l}{\footnotesize Relative FID: ref.\ = 128-step samples, $n=504$.
Absolute FID-50k: ref.\ = real CelebA-HQ 256, $n=50\,000$.}
\end{tabular}
\end{table}

\noindent Despite the difference in absolute values arising from measurement
methodology and sample scale, \textbf{the qualitative trend is fully
consistent}: FID decreases monotonically with increasing NFE, and the largest
quality gap occurs at low NFE (1$\to$4).
This supports the core hypothesis of Shortcut Models: a single network can
deliver flexibly adjustable sample quality within a fixed compute budget.

% ────────────────────────────────────────────────
\subsection{Limitations}
\label{sec:limitations}

\paragraph{1. Dataset scope.}
Experiments were conducted solely on CelebA-HQ 256.
Results were not validated on ImageNet-256 or more diverse datasets, so the
model's generalization to more complex data distributions cannot be assessed.

\paragraph{2. Sample size and relative FID.}
$n = 504$ is $99\times$ smaller than the FID-50k standard.
Relative FID (against 128-step samples) does not measure absolute fidelity
against real images; it only captures within-model quality degradation as NFE
decreases.

\paragraph{3. Seed variance.}
Only 3 seeds were used, making the estimate $\sigma = 0.081$ statistically
unreliable.
Best practice requires at least 5 seeds and bootstrap confidence intervals.

\paragraph{4. OOD baseline is not a true FM model.}
The \texttt{force\_dt0} condition produces a shortcut model operating under
distribution shift, not a separately trained Flow Matching model.
A proper ablation would require a checkpoint trained with
\texttt{--train\_type naive}.

\paragraph{5. Compute infrastructure.}
Running on CPU-only macOS ARM without JIT acceleration
(\texttt{@jax.jit} disabled due to OOM) results in substantially slower
inference compared to the authors' TPU-v3 environment, limiting the feasibility
of scaling to larger sample sizes.

\paragraph{6. No absolute FID.}
The absence of real-image statistics (\texttt{celeba256\_fidstats\_real})
means it is not possible to directly confirm the authors' reported values of
FID-50k $= 6.9$ at 128-step or FID-50k $= 20.5$ at 1-step.
